% 第3章 定义与必要条件
\section{可交易性的最小模型：定义与必要条件}\label{sec:definition}

本章做三件事：给“可交易缝隙”一个能检验的定义；在最小模型里推出它存在的必要条件；把这条条件改写成可以拿数据估的判据。三者合起来，构成后文两个核心命题的共同底座。

\subsection{记号与设定}

把讨论对象限定在“策略”这一抽象层。策略 $\pi$ 是一条从信息到仓位的完整映射：某种信号生成规则给出预测，某种组合规则把它变成权重，某种执行方式把权重变成成交。策略的毛超额收益记 $r_\pi$，摩擦成本（佣金、税费、买卖价差、冲击）记 $C(\pi)$，与所承担风险相称的补偿记 $\vartheta(\pi)$。前两项是工程量，通常有账可查；第三项是定价量，第 \ref{sec:reference} 章会把它显式写成“暴露乘溢价”。本章先把它当作一个待扣除的项接受下来。

关于符号先说清两处安排，免得后文混淆。信息衰减速率用 $\lambda$，它与第 \ref{sec:economics} 章的 Kyle $\lambda$ 同源——价格对订单流越敏感，私有信息的残余就越快被抹平；风险补偿用 $\vartheta(\pi)$；第 \ref{sec:measurement} 章的双幂模型里，毛收益与成本的尺度弹性分别用 $p$ 与 $q$。

\subsection{可交易缝隙的定义}

\begin{definition}[可交易缝隙]\label{def:gap}
设策略 $\pi$ 在给定风险预算下产生毛预期超额收益 $\mathbb{E}[r_\pi]$、交易与冲击成本 $C(\pi)$、以及与所承担风险相称的补偿 $\vartheta(\pi)$。称该策略存在\emph{可交易缝隙}，当且仅当扣除摩擦后的净风险调整收益
\begin{equation}
  G(\pi) \;=\; \mathbb{E}[r_\pi] - C(\pi) - \vartheta(\pi)
  \label{eq:G}
\end{equation}
在\textbf{多重检验校正后}仍然显著大于零，且在样本外与跨样本不反号。
\end{definition}

定义里三个限定词各挡一类错误，缺一个都会漏进假缝隙。

第一，$C(\pi)$ 不是可以事后忽略的尾项。它常常就是决定生死的那一项：Martin 与 Sch\"oneborn（2011）\cite{martin2011} 对均值回归策略的结论一句话可以概括——能赚，但要付成本。把成本放在括号外、当作“实现细节”处理的研究，等于默认了自己要证明的结论。

第二，“多重检验校正后”是因为选择本身的算术：在一个试验过成千上万次的策略空间里，总能挑出一条回测最好的，它的“显著”必须相对于试过多少次而言。Harvey、Liu 与 Zhu（2016）\cite{harveylz2016} 给出的高门槛（$|t|>3$ 一档）与 White（2000）\cite{white2000} 的现实检验，都是这条算术的直接后果。第 \ref{sec:criteria} 章会把这条纪律落成具体的统计流程（DSR、PBO 与试验账本）。

第三，“样本外与跨样本不反号”是最后的诚实检验。一个只在一个样本、一段时期、一个市场上成立的方向，更可能是这段历史的特征而不是市场的特征。第 \ref{sec:reference} 章会把“跨样本”一般化为“跨参照系”，那是这个限定词的理论升格。

定义有意留了一个缺口：它没有规定“在什么观测频率上交易”。这不是疏忽，是下一章的主题——同一个策略在不同的观测网格上，$G(\pi)$ 的每一项都会变，缝隙的存在性因此不能脱离观测装置来陈述。

\subsection{必要条件：信息幅度与成本的赛跑}\label{sec:necessary}

现在把模型降到最简，看看存在一条缝隙至少需要什么。设信号在时刻 $0$ 诞生，其后残余预测幅度按
\begin{equation}
  \delta(t) = \delta_0\,e^{-\lambda t}, \qquad \lambda>0
\end{equation}
衰减，其中 $\delta_0$ 是初始幅度（信号在信息尚新鲜时能兑现的毛收益），$\lambda$ 是信息吸收速率。观测与调仓落在间隔为 $\Delta$ 的网格上，信号从诞生到被交易的期望延迟按半周期计，为 $\Delta/2$；每次调仓付出单次比例成本 $c$。按网格周期分摊，单位时间的净收益为
\begin{equation}
  g(\Delta) \;=\; \frac{\delta_0\,e^{-\lambda\Delta/2} - c}{\Delta}.
  \label{eq:gdelta}
\end{equation}

\begin{proposition}[缝隙存在的必要条件]\label{prop:necessary}
在模型 \eqref{eq:gdelta} 下：存在 $\Delta>0$ 使 $g(\Delta)>0$，当且仅当
\begin{equation}
  \delta_0 > c .
\end{equation}
进一步，可行集合是区间 $(0,\Delta_{\mathrm{hi}})$，其上界
\begin{equation}
  \Delta_{\mathrm{hi}} \;=\; \frac{2}{\lambda}\,\ln\!\frac{\delta_0}{c},
  \label{eq:deltahi}
\end{equation}
且在该区间内 $g$ 严格递减。
\end{proposition}

\begin{proof}
记 $u(\Delta)=\delta_0 e^{-\lambda\Delta/2}$。充分性：若 $\delta_0>c$，由连续性，$\Delta\to 0^{+}$ 时 $u\to\delta_0>c$，故存在小 $\Delta$ 使分子为正、$g>0$。必要性：$g(\Delta)>0$ 蕴含 $u(\Delta)>c$，而 $u(\Delta)\le\delta_0$，故 $\delta_0>c$。上界：$u(\Delta)>c$ 等价于 $\Delta<\frac{2}{\lambda}\ln(\delta_0/c)$。单调性：对 $g$ 求导，
\[
  g'(\Delta)=\frac{-u(\Delta)\bigl(1+\tfrac{\lambda\Delta}{2}\bigr)+c}{\Delta^2},
\]
在可行域 $u>c$ 内，分子小于 $-c(1+\frac{\lambda\Delta}{2})+c<0$，故 $g$ 严格递减。
\end{proof}

命题 \ref{prop:necessary} 把三件事绑在了一起。\textbf{信号的初始幅度必须超过单次成本，否则再好的模型、再聪明的调仓节奏都无力回天。}可行上界 $\Delta_{\mathrm{hi}}$ 对 $\lambda$ 单调递减：\textbf{信息消失得越快，可用的调仓周期就越短。}这解释了为什么高频策略对成本极度敏感——它们的赖以为生的信号按定义衰减极快，可用网格又必须极密，两头夹逼。

但命题还给出一个必须如实交代的副产品：在可行域内 $g$ 严格递减，意味着\textbf{在这个最简模型里，“越密越好”——只要网格还能密下去}。这显然与实务经验相悖：没有人会因此认为无限提高调仓频率是免费的午餐。原因不在推导，而在模型缺了两块结构。第一块是观测与执行的上限：网格不可能无限密，机器与制度都存在最小分辨率，且更密的网格上统计分辨力会下降（第 \ref{sec:measurement} 章）。第二块是成本的凹凸性：单次比例成本 $c$ 为常数是简化，真实的冲击成本对单笔规模呈次线性（平方根），这意味着把交易拆得越碎，单位成交量的冲击越贵（第 \ref{sec:criteria} 章的成本口径）。补上这两块之后，“越密越好”会被一对相反的力取代——这正是下一章要证明的东西。\textbf{最小模型的价值在于指出底线：$\delta_0>c$ 是任何缝隙的必要条件；它的局限在于：光有底线，还画不出形状。}

\subsection{从命题到判据：估计与置信下界}

一条只写着“$\delta_0>c$”的判据是危险的，因为它看起来像同义反复——能赚钱的策略当然幅度盖过成本。这个印象能被打破，条件是两个量都能用可核验的方式估出来，且判据用估计误差修正后的形式。

$c$ 的估计用交易记录直接做：每笔成交均价与决策时点参考价的差额，加上佣金与税费，再除以成交金额，得到该笔的比例成本；多笔取中位数，避免极端单笔带偏。印花税这类制度性费率有公开日历，逐日代入即可；冲击部分若用参与率模型，则按该笔成交名义额与该股 ADV 之比计算（第 \ref{sec:criteria} 章）。

$\delta_0$ 的估计稍微绕一点，用“事件窗平均实现收益”作代理：在信号触发后的 $k$ 期内，按信号方向构建组合的平均超额收益，取事件样本均值 $\hat\delta_0$。它衡量的是信号在信息尚未大面积衰减时能兑现多少。估计量的标准误约为 $\sigma_r/\sqrt{m}$，其中 $m$ 是事件数、$\sigma_r$ 是单事件收益的分散度；$k$ 越大，$\hat\delta_0$ 被衰减压低，故 $k$ 应取信息半衰期的量级。

由于两端都有误差，判据必须用\textbf{置信下界}而不是点估计：
\begin{equation}
  \hat\delta_0 - z_{\alpha}\,\widehat{\mathrm{SE}}(\hat\delta_0) \;>\; \hat c ,
  \label{eq:clb}
\end{equation}
其中 $z_\alpha$ 取单侧分位（$5\%$ 对应约 $1.645$）。这一改动看似保守，实则必要：事件数少、分散度大时，用点估计判断会\textbf{系统性地高估缝隙的存在性}，而这正是大量小样本研究的通病。

\subsection{一个诊断用法}

这套语言有一个立等可用的场景：某策略样本内漂亮、实盘亏钱。此时先按上面的办法估 $\hat\delta_0$ 与 $\hat c$，若置信下界 \eqref{eq:clb} 已经落在 $\hat c$ 之下，问题不在模型过时、不在参数没调好，而在缝隙本身不存在或从未存在。把诊断做在调参之前，可以省下大量对噪声的优化工作。

\subsection{本章要点}

定义给出了缝隙的三重限定（成本、多重检验、样本外不反号），必要条件给出了它的底线（$\delta_0>c$）与时间结构（可行上界 $\Delta_{\mathrm{hi}}$ 随 $\lambda$ 收紧）。判据要可用，必须用置信下界形式。而这一切还悬在一个未说明的假设上——观测网格是给定的。下一章把“观测”本身变成变量。